\section{Discussion, Limitations, and Threats to Validity}
\label{sec:discussion}

\subsection{Threats to Internal and External Validity}
\label{sec:threats_validity}

\noindent\textbf{Circular Stratification Bias:} Our ground-truth strata (S1--S5) were defined using similarity percentiles under the seed vector $\mathbf{v}_0$. Evaluating the cold-start gate on these strata introduces an inherent circularity that tends to favor the seed embedding baseline. We mitigated this threat by applying Horvitz-Thompson sampling weights ($w_h = N_h/n_h$) to reproduce true finite-population distributions, and by sampling low-similarity strata (S3--S5) uniformly at random within their respective percentile bands.

\noindent\textbf{Confounding of Reviewer and Consumer Identity:} In our single-user pilot deployment, the human reviewer (triaging borderline abstracts) and the end user (rating delivered summaries) are the same individual acting at different temporal moments. In institutional deployments with separate screening staff, inter-annotator disagreement could introduce label noise.

\noindent\textbf{Simulated vs. Real Feedback Dynamics:} Offline learning curves (E2--E4) were generated using a deterministic simulated-feedback oracle on labeled development data. While necessary for benchmarking due to sample size constraints, simulated oracles assume static user preferences, whereas real human preferences may exhibit temporal drift or conceptual evolution.

\subsection{System Scalability and Operational Bounds}
\label{sec:scalability}

\noindent\textbf{Percentile Threshold Mechanics:} Using fixed reference percentiles ($P_{\text{high}}=98.0, P_{\text{low}}=90.0$) fixes the expected \textit{proportion} of candidate papers routed to each class rather than an absolute semantic score. If the incoming publication stream undergoes an extreme distribution shift (e.g., a major conference deadline producing an influx of high-quality papers), fixed percentiles may route genuinely relevant manuscripts into overflow or borderline queues.

\noindent\textbf{Queue Overflow Policies:} Dropping overflow papers beyond the 15-item queue cap protects researcher attention but introduces an unavoidable recall penalty during high-volume surges. Future iterations should incorporate adaptive queue sizing based on user review velocity.
